\section{Marco teórico y conceptual}
\label{sec:marco_teorico}

Fundamentación teórica y conceptos clave de la disciplina productiva y del modelo de negocio \parencite{osterwalder2010generacion, kotler2016direccion}.

\subsection{Fundamentos técnicos y herramientas}
\LaTeX{} es un sistema de composición de textos, orientado a la creación de documentos escritos que presentan una alta calidad tipográfica. Por sus características y posibilidades, es ampliamente utilizado en la generación de artículos científicos, libros y tesis \parencite{latexproject}.

\subsection{Imágenes y Gráficos}
Para incluir figuras dentro de la documentación, se recomienda ubicar las imágenes en la carpeta \texttt{imagenes/} y hacer referencia a ellas utilizando el entorno \texttt{figure}.

\begin{figure}[htbp]
    \centering
    % Ejemplo de inclusión de imagen:
    % \includegraphics[width=0.6\textwidth]{ejemplo.png}
    \rule{0.6\textwidth}{4cm} % Marco temporal tipo placeholder
    \caption{Representación gráfica esquemática del modelo productivo.}
    \label{fig:esquema_ejemplo}
\end{figure}

Como se observa en la Figura \ref{fig:esquema_ejemplo}, las representaciones visuales permiten clarificar la arquitectura del sistema y los flujos de trabajo planteados.
